% ============================================================================
% MATEMÁTICAS 2 · ESTILO BEAMER MODERNO Y MINIMALISTA
% ============================================================================
% Diseño: fondo blanco, títulos a la izquierda, acento azul discreto,
% viñetas claras, márgenes amplios y número de diapositiva pequeño.
%
% Compatible con los comandos usados en los temas:
%   \dd, \figura, \portadatema, \portadaasignatura.
% No requiere temas ni fuentes externas.
% ============================================================================

% ----------------------------------------------------------------------------
% 1. IDIOMA, TIPOGRAFÍA Y PAQUETES
% ----------------------------------------------------------------------------
\usepackage[T1]{fontenc}
\usepackage[utf8]{inputenc}
\usepackage[spanish,provide=*]{babel}
\usepackage{lmodern}
\usepackage{bm}
\usepackage{amsmath,amssymb,mathtools,graphicx,array}
\usepackage{tikz}
\usetikzlibrary{calc}
\newcolumntype{C}[1]{>{\centering\arraybackslash}p{#1}}
\usepackage[most]{tcolorbox}
% ============================================================================
% DIAPOSITIVA DE CONTEXTO: CUATRO APLICACIONES
% Incluir en el preámbulo (config.tex), una sola vez:
%   \usepackage[most]{tcolorbox}
%   % ============================================================================
% DIAPOSITIVA DE CONTEXTO: CUATRO APLICACIONES
% Incluir en el preámbulo (config.tex), una sola vez:
%   \usepackage[most]{tcolorbox}
%   % ============================================================================
% DIAPOSITIVA DE CONTEXTO: CUATRO APLICACIONES
% Incluir en el preámbulo (config.tex), una sola vez:
%   \usepackage[most]{tcolorbox}
%   \input{contexto_estilos.tex}
% Después, pegar el frame que aparece al final del archivo en tema1.tex.
% ============================================================================

% --- Colores de las tarjetas: fondos suaves + acentos contrastados ---
\definecolor{ctxAmarillo}{HTML}{C18A25}
\definecolor{ctxAmarilloFondo}{HTML}{FFF5DD}
\definecolor{ctxRojo}{HTML}{C25B58}
\definecolor{ctxRojoFondo}{HTML}{FCEBE9}
\definecolor{ctxAzul}{HTML}{397BB6}
\definecolor{ctxAzulFondo}{HTML}{EAF3FB}
\definecolor{ctxVerde}{HTML}{3C8A63}
\definecolor{ctxVerdeFondo}{HTML}{EAF5ED}

% Placeholder de imagen (sustituible por \includegraphics).
% Argumentos: color de borde, descripción breve.
\newcommand{\imagenpendiente}[2]{%
  \fcolorbox{#1!50!white}{white}{%
    \parbox[c][1.75cm][c]{1.72cm}{%
      \centering\color{mateMuted}\fontsize{8}{9}\selectfont #2%
    }%
  }%
}

% Tarjeta de aplicación: color, fondo, imagen, título, explicación, fórmula.
\newcommand{\tarjetacontexto}[6]{%
  \begin{tcolorbox}[
    enhanced,
    width=\linewidth,
    height=3.12cm,
    colback=#2,
    colframe=#1!48!white,
    boxrule=0.35pt,
    arc=1.7mm,
    boxsep=0pt,
    left=2.2mm,right=1.8mm,top=1.5mm,bottom=1.5mm,
    borderline west={1.4mm}{0pt}{#1},
    valign=center
  ]
    \begin{minipage}[c]{0.30\linewidth}
      \centering #3
    \end{minipage}%
    \hfill
    \begin{minipage}[c]{0.67\linewidth}
      \raggedright
      \parbox[t][0.93cm][t]{\linewidth}{%
        \raggedright\color{#1!70!black}\bfseries\fontsize{10.7}{12.2}\selectfont #4%
      }\par
      \parbox[t][1.30cm][t]{\linewidth}{%
        \raggedright\color{mateText}\fontsize{9.1}{11.3}\selectfont #5%
      }\par
      \parbox[t][0.55cm][t]{\linewidth}{%
        \color{mateText}\fontsize{11.2}{13}\selectfont \(\displaystyle #6\)%
      }%
    \end{minipage}
  \end{tcolorbox}%
}


% Después, pegar el frame que aparece al final del archivo en tema1.tex.
% ============================================================================

% --- Colores de las tarjetas: fondos suaves + acentos contrastados ---
\definecolor{ctxAmarillo}{HTML}{C18A25}
\definecolor{ctxAmarilloFondo}{HTML}{FFF5DD}
\definecolor{ctxRojo}{HTML}{C25B58}
\definecolor{ctxRojoFondo}{HTML}{FCEBE9}
\definecolor{ctxAzul}{HTML}{397BB6}
\definecolor{ctxAzulFondo}{HTML}{EAF3FB}
\definecolor{ctxVerde}{HTML}{3C8A63}
\definecolor{ctxVerdeFondo}{HTML}{EAF5ED}

% Placeholder de imagen (sustituible por \includegraphics).
% Argumentos: color de borde, descripción breve.
\newcommand{\imagenpendiente}[2]{%
  \fcolorbox{#1!50!white}{white}{%
    \parbox[c][1.75cm][c]{1.72cm}{%
      \centering\color{mateMuted}\fontsize{8}{9}\selectfont #2%
    }%
  }%
}

% Tarjeta de aplicación: color, fondo, imagen, título, explicación, fórmula.
\newcommand{\tarjetacontexto}[6]{%
  \begin{tcolorbox}[
    enhanced,
    width=\linewidth,
    height=3.12cm,
    colback=#2,
    colframe=#1!48!white,
    boxrule=0.35pt,
    arc=1.7mm,
    boxsep=0pt,
    left=2.2mm,right=1.8mm,top=1.5mm,bottom=1.5mm,
    borderline west={1.4mm}{0pt}{#1},
    valign=center
  ]
    \begin{minipage}[c]{0.30\linewidth}
      \centering #3
    \end{minipage}%
    \hfill
    \begin{minipage}[c]{0.67\linewidth}
      \raggedright
      \parbox[t][0.93cm][t]{\linewidth}{%
        \raggedright\color{#1!70!black}\bfseries\fontsize{10.7}{12.2}\selectfont #4%
      }\par
      \parbox[t][1.30cm][t]{\linewidth}{%
        \raggedright\color{mateText}\fontsize{9.1}{11.3}\selectfont #5%
      }\par
      \parbox[t][0.55cm][t]{\linewidth}{%
        \color{mateText}\fontsize{11.2}{13}\selectfont \(\displaystyle #6\)%
      }%
    \end{minipage}
  \end{tcolorbox}%
}


% Después, pegar el frame que aparece al final del archivo en tema1.tex.
% ============================================================================

% --- Colores de las tarjetas: fondos suaves + acentos contrastados ---
\definecolor{ctxAmarillo}{HTML}{C18A25}
\definecolor{ctxAmarilloFondo}{HTML}{FFF5DD}
\definecolor{ctxRojo}{HTML}{C25B58}
\definecolor{ctxRojoFondo}{HTML}{FCEBE9}
\definecolor{ctxAzul}{HTML}{397BB6}
\definecolor{ctxAzulFondo}{HTML}{EAF3FB}
\definecolor{ctxVerde}{HTML}{3C8A63}
\definecolor{ctxVerdeFondo}{HTML}{EAF5ED}

% Placeholder de imagen (sustituible por \includegraphics).
% Argumentos: color de borde, descripción breve.
\newcommand{\imagenpendiente}[2]{%
  \fcolorbox{#1!50!white}{white}{%
    \parbox[c][1.75cm][c]{1.72cm}{%
      \centering\color{mateMuted}\fontsize{8}{9}\selectfont #2%
    }%
  }%
}

% Tarjeta de aplicación: color, fondo, imagen, título, explicación, fórmula.
\newcommand{\tarjetacontexto}[6]{%
  \begin{tcolorbox}[
    enhanced,
    width=\linewidth,
    height=3.12cm,
    colback=#2,
    colframe=#1!48!white,
    boxrule=0.35pt,
    arc=1.7mm,
    boxsep=0pt,
    left=2.2mm,right=1.8mm,top=1.5mm,bottom=1.5mm,
    borderline west={1.4mm}{0pt}{#1},
    valign=center
  ]
    \begin{minipage}[c]{0.30\linewidth}
      \centering #3
    \end{minipage}%
    \hfill
    \begin{minipage}[c]{0.67\linewidth}
      \raggedright
      \parbox[t][0.93cm][t]{\linewidth}{%
        \raggedright\color{#1!70!black}\bfseries\fontsize{10.7}{12.2}\selectfont #4%
      }\par
      \parbox[t][1.30cm][t]{\linewidth}{%
        \raggedright\color{mateText}\fontsize{9.1}{11.3}\selectfont #5%
      }\par
      \parbox[t][0.55cm][t]{\linewidth}{%
        \color{mateText}\fontsize{11.2}{13}\selectfont \(\displaystyle #6\)%
      }%
    \end{minipage}
  \end{tcolorbox}%
}



% Notación matemática.
\newcommand{\dd}{\mathop{}\!\mathrm{d}}

% ----------------------------------------------------------------------------
% 2. TEMA BASE Y PALETA DE COLORES
% ----------------------------------------------------------------------------
\usetheme{default}
\usefonttheme{professionalfonts}

% Ajustar estos cuatro colores basta para cambiar la identidad visual.
\definecolor{mateBlueDark}{RGB}{31,75,108}   % títulos y elementos activos
\definecolor{mateText}{RGB}{36,44,52}        % texto principal
\definecolor{mateMuted}{RGB}{104,117,128}    % información secundaria
\definecolor{mateRule}{RGB}{218,225,230}     % separadores muy discretos
\definecolor{mateSoft}{RGB}{246,248,250}     % fondo de bloques opcionales

\setbeamercolor{background canvas}{bg=white}
\setbeamercolor{normal text}{fg=mateText,bg=white}
\setbeamercolor{frametitle}{fg=mateBlueDark,bg=white}
\setbeamercolor{structure}{fg=mateBlueDark}
\setbeamercolor{alerted text}{fg=mateBlueDark}
\setbeamercolor{itemize item}{fg=mateBlueDark}
\setbeamercolor{itemize subitem}{fg=mateBlueDark}
\setbeamercolor{block title}{fg=mateBlueDark,bg=mateSoft}
\setbeamercolor{block body}{fg=mateText,bg=mateSoft}

% Márgenes del área de contenido (el título y el pie usan los mismos márgenes).
\setbeamersize{text margin left=8mm,text margin right=8mm}

% Sin iconos de navegación, sombras, fondos degradados ni decoraciones.
\setbeamertemplate{navigation symbols}{}
\setbeamercovered{invisible}
\setbeamertemplate{background canvas}[default]

% ----------------------------------------------------------------------------
% 3. TIPOGRAFÍA Y VIÑETAS
% ----------------------------------------------------------------------------
\setbeamerfont{frametitle}{family=\sffamily,series=\bfseries,size=\Large}
\setbeamerfont{footline}{family=\sffamily,size=\scriptsize}

% Puntos discretos en lugar de iconos ornamentales.
\setbeamertemplate{itemize item}[circle]
\setbeamertemplate{itemize subitem}[circle]

% Párrafos separados de forma natural, sin \vspace generalizado.
\setlength{\parindent}{0pt}
\setlength{\parskip}{0em}

% ----------------------------------------------------------------------------
% 4. TÍTULO DE CADA DIAPOSITIVA
% ----------------------------------------------------------------------------
% Alineado a la izquierda, sin cabecera coloreada.
% Una línea gris tenue separa el título del contenido.
\setbeamertemplate{frametitle}{%
  \nointerlineskip%
  \vspace*{0.20cm}%
  \begin{beamercolorbox}[
    wd=\paperwidth,
    leftskip=8mm,
    rightskip=8mm,
    sep=0pt
  ]{frametitle}%
    \usebeamerfont{frametitle}\insertframetitle\par%
    \vspace*{0.12cm}%
  \end{beamercolorbox}%
  \nointerlineskip%
  \textcolor{mateRule}{%
    \rule{\dimexpr\paperwidth-16mm\relax}{0.45pt}%
  }\par%
  \vspace*{0.07cm}%
}

% ----------------------------------------------------------------------------
% 5. PIE DE PÁGINA
% ----------------------------------------------------------------------------
% Etiqueta discreta a la izquierda y número de transparencia a la derecha.
\setbeamercolor{footline}{fg=mateMuted,bg=white}
\setbeamertemplate{footline}{%
  \leavevmode%
  \begin{beamercolorbox}[
    wd=\paperwidth,
    ht=0.25cm,
    dp=0.08cm,
    leftskip=8mm,
    rightskip=8mm
  ]{footline}%
    \usebeamerfont{footline}%
    \hfill\insertframenumber%
  \end{beamercolorbox}%
}

% ----------------------------------------------------------------------------
% 6. DATOS DE LA ASIGNATURA
% ----------------------------------------------------------------------------
\title{Matemáticas 2}
\subtitle{Grado en Ingeniería Informática}
\author{}
\institute{Universidad de Alicante}
\date{}

% ----------------------------------------------------------------------------
% 7. PORTADA AUTOMÁTICA DE CADA TEMA
% ----------------------------------------------------------------------------
% Se inserta al escribir \section{Nombre del tema}.
% No hay formas flotantes, ni degradados, ni cuadros ornamentales.
\AtBeginSection[]{%
  \begin{frame}[plain]
    \vspace*{0.9cm}
    {\color{mateMuted}\large\bfseries
      TEMA \insertsectionnumber\par}
    \vspace{0.40cm}
    {\color{mateBlueDark}\bfseries
      \fontsize{29}{35}\selectfont\insertsectionhead\par}
    \vspace{0.45cm}
    {\color{mateBlueDark}\rule{1.5cm}{1.6pt}\par}
    \vfill
    {\color{mateMuted}\small
      Matemáticas 2 · Grado en Ingeniería Informática\par}
    \vspace*{0.25cm}
  \end{frame}%
}

% ----------------------------------------------------------------------------
% 8. PORTADA AUTOMÁTICA DE CADA SUBSECCIÓN
% ----------------------------------------------------------------------------
% Se inserta al escribir \subsection{Nombre del apartado}.
\AtBeginSubsection[]{%
  \begin{frame}[plain]
    \vspace*{1.20cm}
    {\color{mateMuted}\normalsize
      Tema \insertsectionnumber · \insertsectionhead\par}
    \vspace{0.35cm}
    {\color{mateBlueDark}\bfseries
      \fontsize{25}{31}\selectfont\insertsubsectionhead\par}
    \vspace{0.35cm}
    {\color{mateRule}\rule{\linewidth}{0.5pt}\par}
    \vfill
  \end{frame}%
}

% ----------------------------------------------------------------------------
% 9. COMANDOS AUXILIARES EMPLEADOS EN LOS TEMAS
% ----------------------------------------------------------------------------

% Figura centrada, con anchura y altura máximas relativas al texto.
% Uso: \figura{archivo}{ancho relativo}{alto relativo}
\newlength{\mateFigureHeight}
\newcommand{\figura}[3]{%
  \begingroup
  \setlength{\mateFigureHeight}{#3\textheight}%
  \ifdim\mateFigureHeight>.80\textheight
    \setlength{\mateFigureHeight}{.80\textheight}%
  \fi
  \begin{center}
    \includegraphics[width=#2\textwidth,height=\mateFigureHeight,keepaspectratio]{#1}%
  \end{center}%
  \endgroup
}

% Mantiene la numeración de los temas cuando se compilan individualmente.
% Uso: \portadatema{1}{Funciones de varias variables}
\newcommand{\portadatema}[2]{%
  \setcounter{section}{\numexpr#1-1\relax}%
  \section{#2}%
}

% Portada opcional de la asignatura (desde main.tex).
\newcommand{\portadaasignatura}{%
  \begin{frame}[plain]
    \vspace*{1.30cm}
    {\color{mateBlueDark}\fontsize{34}{40}\selectfont\bfseries
      Matemáticas 2\par}
    \vspace{0.35cm}
    {\color{mateText}\Large Grado en Ingeniería Informática\par}
    \vfill
    {\color{mateMuted}\small Universidad de Alicante\par}
    \vspace*{0.25cm}
  \end{frame}%
}
